\begin{helper}
In the actual split-outside-blocker setting, assume the two
activation--priority laws have the Bernoulli activation atom formulas and the
finite lower-rectangle priority formulas with
\(0\le\gamma/\Delta\le 1\).  Let \(\phi:V\to V'\) be injective, let
\(r\in V\), let \(X\subseteq N_G(r)\), and suppose that \(\phi\) preserves
all adjacencies inside \(X\cup\{r\}\).  Let
\[
  \beta(x,z)\in V'
\]
be the private replacement vertex attached to each outside incidence
\(x\in X\), \(z\notin X\cup\{r\}\), \(xz\in E(G)\).  Assume that \(\beta\) is
injective, fresh from \(\phi(X\cup\{r\})\), and adjacent to the intended
vertex \(\phi(x)\).  Also assume the full local replacement conditions:
\(\beta(x,z)\) is adjacent to no embedded candidate \(\phi(y)\) with
\(y\in X\) except when \(y=x\), it is nonadjacent to \(\phi(r)\), and every
neighbor of \(\phi(x)\) in \(G'\), for \(x\in X\), is either an embedded
neighbor \(\phi(y)\) with \(y\in X\cup\{r\}\) and \(xy\in E(G)\), or is one of
the private vertices \(\beta(x,z)\) attached to an outside source neighbor
\(z\notin X\cup\{r\}\) of \(x\).

Then there is a single actual witness surface consisting of an injective
enumeration \(z_0,\ldots,z_{n-1}\) of all outside blockers adjacent to
\(X\), candidate sets \(B_k=\{x\in X:xz_k\in E(G)\}\), source and target
coordinate layers, a finite atom index set \(\Omega\), paired cells
\(C_\omega,C'_\omega\), and common weights \(w(\omega)\).  These witnesses
satisfy all coordinate clauses of
\(\noderef{SamplingSplitOutsideBlockerActualCoordinatePackage}\) and all
prescribed common refined-cell clauses of
\(\noderef{SamplingSplitOutsideBlockerPrescribedCommonRefinedCellPackage}\):
cell masses agree with \(w\) on both sides, source and target cells are
pairwise disjoint, they cover the unit-priority supports up to the stated null
complements, and cell membership is transported by equality of the common
activation trace and common priority coordinates.  In addition, each returned
source cell \(C_\omega\) and target cell \(C'_\omega\) is measurable, and the
source endpoint event \(I_G(A,\pi)\cap X\ne\emptyset\) and target endpoint
event \(I_{G'}(A',\pi')\cap\phi(X)\ne\emptyset\) are measurable in the
corresponding activation--priority product spaces.

On this same \(n,z,B,\Omega,C,C',w\) surface there is a
\(\noderef{SamplingSplitOutsideBlockerActualSameSurfaceInputData}\) package.
Equivalently, one has a paired mixed measure surface together with the source
endpoint identity, the target endpoint identity, and the adjacent
shared-context decompositions for the consecutive mixed masses.  The theorem
does not assert a same-index source--target stage equality and does not sum
paired outside-\(B_k\) primitive atom equalities; it records the concrete
same-surface input that must be built below the hybrid mass law.
\end{helper}
\begin{proof}
Apply
\(\noderef{SamplingSplitOutsideBlockerActualCoordinatePackage}\) to
\(G,G',\phi,r,X\), and \(\beta\), using the injectivity, freshness, and
intended-adjacency hypotheses required by that coordinate package.  Its
statement returns the actual
outside-blocker list \(z_0,\ldots,z_{n-1}\), the sets \(B_k\), the source
common and extension coordinates, the target common and extension
coordinates, and the full target coordinate set.  The returned clauses give
the injectivity of \(k\mapsto z_k\), the fact that every listed \(z_k\) lies
outside \(X\cup\{r\}\) and is adjacent to some vertex of \(X\), the converse
coverage of every outside blocker adjacent to \(X\), the identity
\(B_k=\{x\in X:xz_k\in E(G)\}\), the source-coordinate disjointness and
containment clauses, the identity of the target common coordinates with the
\(\phi\)-image of the source common coordinates, the target-coordinate
disjointness and union clauses, and the containment of every relevant
private blocker \(\beta(x,z_k)\) in the target extension coordinates.

Now keep exactly these witnesses fixed and apply
\(\noderef{SamplingSplitOutsideBlockerPrescribedCommonRefinedCellPackage}\).
Because that node is prescribed on the supplied coordinate surface, it does
not change \(n,z_k,B_k\) or any coordinate set.  It returns a finite atom type
\(\Omega\), source cells \(C_\omega\), target cells \(C'_\omega\), and weights
\(w(\omega)\).  Its statement supplies \(w(\omega)\ge0\), the two mass
identities
\[
  \operatorname{ofReal}(w(\omega))=\nu(C_\omega),
  \qquad
  \operatorname{ofReal}(w(\omega))=\nu'(C'_\omega),
\]
source and target pairwise disjointness, source and target unit-support
cover, the two null-complement identities, and the transported
common-coordinate membership equivalence for each \(\omega\).

It remains to build the actual same-surface input on this very atom surface.
First apply
\(\noderef{SamplingSplitOutsideBlockerMarginalPreservingCouplingConstruction}\)
to the same \(\nu,\nu',\phi,\texttt{sourceCoord},
\texttt{targetCommonCoord},\Omega,C,C'\), and \(w\).  Its hypotheses are
the activation and priority product laws from the present theorem, the
identity
\(\texttt{targetCommonCoord}=\phi(\texttt{sourceCoord})\) from the coordinate
package, and the nonnegativity, mass, disjointness, transported-membership,
and measurability clauses from the prescribed common-refined-cell package.
It gives paired cells
\(\widehat C_\omega\subseteq C_\omega\times C'_\omega\) and one paired
measure \(\lambda\) on the same atom index set \(\Omega\), with
\[
  \lambda(\widehat C_\omega)=\operatorname{ofReal}(w(\omega)).
\]
Moreover, every paired state in \(\widehat C_\omega\) synchronizes the
common activation trace and common priority coordinates, and the restriction
of \(\lambda\) to \(\widehat C_\omega\) has source marginal
\(\nu|_{C_\omega}\) and target marginal \(\nu'|_{C'_\omega}\).  This is the
registered paired-coupling layer; no endpoint or adjacent comparison is
obtained from total mass equality alone.
The strengthened constructor also exports the exact paired-cell formula
and a \noderef{SamplingSplitOutsideBlockerProductRealization} for this
same $\lambda$. In particular it is the push-forward of the independent
Bernoulli--uniform tag law, not an arbitrary coupling with those marginals.
Apply \noderef{SamplingSplitOutsideBlockerProductAtomContextInvariant}
using its tag-overlap clause, the prescribed membership transport, and the
exact paired-cell formula. The inverse paired cell is a source-cell
pullback and is unchanged whenever all common-tag coordinates are fixed.
Each source-current tag is noncommon: $z_k$ is in the extension set
disjoint from the common set, and source tags are injective. Each target
private tag is noncommon too: an equality with a common source tag would,
by the overlap clause, put that private vertex in the target common set,
contradicting target extension disjointness. Thus restricting to the paired
atom does not constrain the current tags. This conclusion holds pointwise
on the tagged product space and needs no conditioning on priority values.

Apply
\(\noderef{SamplingSplitOutsideBlockerActualMixedStageMassNormalization}\)
directly to these fixed witnesses and this coupling, using the already proved
nonnegativity of \(w\).  It supplies the displayed paired mixed events
\(\widehat S_j(\omega)\): a single candidate is active on both copies,
passes both embedded-neighbor tests, passes the private tests for \(i<j\),
and passes the retained source tests for \(j\le i<n\).  It also supplies one
family \(M_j(\omega)\), nonnegative for every \(j,\omega\), zero whenever
\(w(\omega)=0\), and satisfying
\[
 \operatorname{ofReal}(w(\omega)M_j(\omega))
 =\lambda(\widehat S_j(\omega)\cap\widehat C_\omega)
 \qquad(j\le n).
\]
Together with the coordinate fields, synchronization, projections and cell
masses supplied by the same coupling, these are all the fields of
\noderef{SamplingSplitOutsideBlockerPairedMixedMeasureSurface}.
Thus the slim carrier is now assembled on these witnesses; it has no
intermediate boundary projections or assumed adjacent inequalities.

For each \(k<n\), use
\noderef{SamplingSplitOutsideBlockerActualAdjacentSharedContextProducer}
on this carrier. It supplies the measurable paired context \(C_x\), the
unchanged success event \(U\) witnessed outside \(B_k\), and the two paired
boundaries \(L,R\), both excluding \(U\). It proves the disjoint identities
\(\widehat S_k=U\mathbin{\dot\cup}L\) and
\(\widehat S_{k+1}=U\mathbin{\dot\cup}R\).
These are identities for paired events, not one-sided pullbacks.
The candidate and all its noncurrent tests use the same witness.

Moreover,
\noderef{SamplingSplitOutsideBlockerPairedContextCurrentInvariant}
applies with injectivity of \(z_i\), injectivity of \(\beta\), the outside
condition, and private-coordinate freshness from the present hypotheses.
It shows that each \(C_x\) is unaffected by changing the source-current
coordinate or any target-current private coordinate. Consequently \(U\)
and the survivor set \(K=\{x\in B_k:C_x\}\) are unchanged under those
coordinate changes. This is an event-level statement only.

Now apply \noderef{SamplingSplitOutsideBlockerActualProductAdjacentIntegral}
on this same carrier, with $p=\gamma/\Delta$, the retained product
realization, and the exact cell formula from the coupling constructor.
Its measurability, weight nonnegativity, and membership-transport
hypotheses are precisely the clauses of the prescribed atom package.
Its remaining coordinate hypotheses follow from disjointness:
$z_i$ is in the source extension set and therefore is not in the source
common set; each $\beta(x,z_i)$ is in the target extension set and
therefore is not in the target common set. Injectivity of $\beta$ is
an original hypothesis. The integral theorem yields, for every $k<n$
and every $\omega$, nonnegative $u,l,r'$ representing the actual paired
piece masses divided by the atom weight when that weight is positive,
and zero when the weight is zero. It proves
\[
 M_k(\omega)=u+l,\qquad M_{k+1}(\omega)=u+r',\qquad l\le r'.
\]
These are exactly the required numerical adjacent certificates. In
particular the theorem uses the same global $M$ supplied by normalization;
no representatives are chosen independently at the two ends of each step.
The fiber calculation and integration used to establish this certificate
are registered below the integral theorem, not assumed as fields of the
carrier.

For the source endpoint apply
\noderef{SamplingSplitOutsideBlockerActualSourceEndpointFromCoupling}
to this carrier and coupling. The paired measure and paired cells agree
by their definitions. Injectivity of $\phi$ is assumed; the outside-blocker
coverage is supplied by the coordinate package; the parameter bounds and
source product laws are original hypotheses. The exact source endpoint
event is measurable by the prescribed atom package. Thus it gives
$\operatorname{ofReal}(wM_0)=\nu(\{I_G\cap X\ne\varnothing\}\cap C_\omega)$.
For the target endpoint apply
\noderef{SamplingSplitOutsideBlockerActualTargetEndpointFromCoupling}.
Again the measures and cells agree definitionally, coverage comes from
the coordinate package, and the product laws, parameter bounds, embedded
adjacency equivalence, intended private adjacency, uniqueness of its
candidate neighbor, nonadjacency to the root, and exhaustive local
neighbor description are all original hypotheses. The target endpoint
measurability is also a conclusion of the atom package. The result is
$\operatorname{ofReal}(wM_n)=\nu'(\{I_{G'}\cap\phi(X)\ne\varnothing\}
\cap C'_\omega)$. These endpoint theorems include their displayed-event
measurability and strict-boundary reduction; no intermediate event is
projected to one side.

The surface, these two endpoint identities, and the adjacent certificates
are precisely the four fields of
\noderef{SamplingSplitOutsideBlockerActualSameSurfaceInputData}.
Assemble them, retaining the original enumeration, coordinate sets, cells,
weights, disjointness, support-cover, null-complement, membership-transport,
and measurability clauses already obtained above. This proves every
conjunct of the requested existential on one and the same witness surface.
\end{proof}
